\section{Motivation}
\label{sec:motivation}

% SKELETON. The structure, the figure and the caption are in place; the prose of each
% paragraph is still to be written. The bracketed notes say what each paragraph should
% argue, and which measured numbers from Section~\ref{sec:evaluation} back it up.
% The numbers quoted in the notes are the ones that section currently reports.

% Three topologies of a migration between two shard groups of 2f+1 = 5 replicas.
% \topopanel{1} no replication at the recipient, {2} migration layered on replication,
% {3} the integrated chain. \topolegend draws the key shared by the three panels.
%
% Each panel is drawn at 1pt per unit and then scaled by \topozoom so that its fixed
% 178pt bounding box lands on the 0.32\textwidth of a subfigure. The scaling is uniform
% and shared by the legend, so the three panels and the key carry one type size.
\definecolor{ink}{RGB}{34,34,34}
\definecolor{irfill}{gray}{0.86}
\newcommand{\topozoom}{0.927}
\tikzset{topostyle/.style={%
    font=\fontsize{7.5}{9}\fontfamily{phv}\selectfont,
    >={Stealth[length=4pt,width=3.2pt]},
    shard/.style={draw=ink,line width=0.6pt,fill=white,rounded corners=1.5pt,
                  minimum width=26pt,minimum height=14pt,inner sep=0pt},
    leader/.style={shard,fill=ink,text=white,font=\fontsize{7.5}{9}\fontfamily{phv}\bfseries},
    bulk/.style={->,draw=ink,line width=1.8pt,>={Stealth[length=5pt,width=5pt]}},
    abulk/.style={bulk,dash pattern=on 3pt off 2pt},
    notify/.style={->,draw=ink,line width=0.6pt,dash pattern=on 2.2pt off 1.4pt,
                   rounded corners=2pt},
    lbl/.style={inner sep=1pt,text=ink},
    phase/.style={rounded corners=2pt},
    rolelbl/.style={lbl,font=\fontsize{7}{8}\fontfamily{phv}\itshape},
    chunk/.style={draw=ink,line width=0.6pt,fill=irfill,minimum width=18pt,
                  minimum height=10pt,inner sep=0pt,
                  font=\fontsize{6.6}{7.4}\fontfamily{phv}\selectfont},
}}

\newcommand{\topopanel}[1]{%
\scalebox{\topozoom}{%
\begin{tikzpicture}[x=1pt,y=-1pt,topostyle]
  % --- geometry --------------------------------------------------------------
  \def\Dspine{6}   \def\Dx{46}           % donor: log spine, replica column
  \def\Rx{136}     \def\Rspine{174}      % recipient: replica column, log spine
  \def\rowa{22} \def\rowb{48} \def\rowc{74} \def\rowd{100} \def\rowe{126}

  % one bounding box for all three panels, so that they scale alike
  \useasboundingbox (0,12) rectangle (178,166);

  % --- the replicas that make the group durable -----------------------------
  \ifnum#1>1 \fill[phase,irfill] (118,67) rectangle (154,133); \fi

  % --- donor replicas --------------------------------------------------------
  \node[shard]  at (\Dx,\rowa) {A1};
  \node[shard]  at (\Dx,\rowb) {A2};
  \node[shard]  at (\Dx,\rowc) {A3};
  \node[shard]  at (\Dx,\rowd) {A4};
  \node[leader] at (\Dx,\rowe) {A5};
  \node[rolelbl] at (\Dx,141) {leader};

  % --- recipient replicas ----------------------------------------------------
  \node[shard]  at (\Rx,\rowa) {B1};
  \node[shard]  at (\Rx,\rowb) {B2};
  \node[shard]  at (\Rx,\rowc) {B3};
  \node[shard]  at (\Rx,\rowd) {B4};
  \node[leader] at (\Rx,\rowe) {B5};
  \node[rolelbl] at (\Rx,141) {leader};

  % --- replication paths -----------------------------------------------------
  \foreach \y in {\rowa,\rowb,\rowc,\rowd} {
    \draw[notify] (\Dx-13,\rowe) -- (\Dspine,\rowe) -- (\Dspine,\y) -- (\Dx-13,\y);
    \draw[notify] (\Rx+13,\rowe) -- (\Rspine,\rowe) -- (\Rspine,\y) -- (\Rx+13,\y);
  }

  % --- leader-to-leader bulk transfer ---------------------------------------
  \draw[bulk] (\Dx+13,\rowe) -- (\Rx-14,\rowe);
  % The layered panel carries the group as one block. The chain panel carries the
  % same block split in two, so the two spans match end to end.
  \ifnum#1=2
    \node[chunk,minimum width=41pt] at (88,118) {\textit{B}};
  \else
    \node[chunk] at (76.5,118) {\textit{B}/2};
    \node[chunk] at (99.5,118) {\textit{B}/2};
  \fi
  % the brace counts the chunks the group is split into
  \draw[ink,line width=0.5pt,decorate,decoration={brace,amplitude=3pt}]
    (67.5,111) -- (108.5,111);
  \ifnum#1=2
    \node[lbl,anchor=base] at (88,104) {$k = 1$};
  \else
    \node[lbl,anchor=base] at (88,104) {$k = 2$};
  \fi
  \node[lbl,inner sep=1.5pt,
        font=\fontsize{7.5}{9}\fontfamily{phv}\bfseries] at (90,137) {R};

  % --- how the data reaches the followers ------------------------------------
  \ifnum#1=2
    % the leader sends every chunk itself, once per follower
    \foreach \y in {\rowa,\rowb,\rowc,\rowd}
      \draw[bulk,rounded corners=2pt] (\Rx-13,122) -- (114,122) -- (114,\y) -- (\Rx-13,\y);
    \node[lbl,anchor=east,font=\fontsize{7.5}{9}\fontfamily{phv}\bfseries] at (110,61) {R};
  \fi
  % --- what the transfer costs ----------------------------------------------
  \ifnum#1=2
    \node[lbl,anchor=base] at (89,162)
      {$n = 5$:\quad $T_{a} = \tfrac{n+1}{2}\,B/R = 3\,B/R$};
  \else
    \node[lbl,anchor=base] at (89,162)
      {$n = 5$:\quad $T_{b} = \bigl(1 + \tfrac{n-1}{2k}\bigr)B/R = 2\,B/R$};
  \fi

  \ifnum#1=3
    % the chain runs up from the leader, the first two hops reach a majority
    \draw[bulk]  (\Rx,118.5) -- (\Rx,107.5);
    \draw[bulk]  (\Rx,92.5)  -- (\Rx,81.5);
    \node[lbl,anchor=east,font=\fontsize{7.5}{9}\fontfamily{phv}\bfseries] at (120,87) {R};
    \draw[abulk] (\Rx,66.5)  -- (\Rx,55.5);
    \draw[abulk] (\Rx,40.5)  -- (\Rx,29.5);
  \fi
\end{tikzpicture}}%
}

% An empty slot that reserves exactly the space of a panel, for a third figure.
\newcommand{\topoblank}{%
\scalebox{\topozoom}{%
\begin{tikzpicture}[x=1pt,y=-1pt,topostyle]
  \useasboundingbox (0,12) rectangle (178,166);
\end{tikzpicture}}%
}

% The model the two panels above describe, drawn: the time to get one group onto a
% majority of the recipient group, against the number of chunks k, for both ways.
% The formula, not measured data.
%   (a) layered, replication follows the transfer, so splitting does not help:
%       (n+1)/2 B/R, and even a pipelined layered way could not pass (n-1)/2 B/R,
%   (b) chain, every link carries B once over (n+1)/2 pipelined hops:
%       (1 + (n-1)/2k) B/R,
%       so it falls towards B/R.
% The y axis title is turned with \rotatebox rather than with a node rotation, because
% y runs downwards here and a node rotation would come out mirrored.
\newcommand{\topoplot}{%
\scalebox{\topozoom}{%
\begin{tikzpicture}[x=1pt,y=-1pt,topostyle,
    axis/.style={draw=ink,line width=0.7pt},
    ticklbl/.style={inner sep=1.5pt,text=ink,
                    font=\fontsize{6.6}{7.4}\fontfamily{phv}\selectfont},
    layered/.style={draw=ink,line width=1.6pt,dash pattern=on 4pt off 2.5pt},
    chain/.style={draw=ink,line width=1.6pt,line join=round},
    grid/.style={draw=ink!45,line width=0.5pt,dash pattern=on 1.2pt off 1.6pt}]
  \useasboundingbox (0,12) rectangle (178,166);

  \node[ticklbl,inner sep=0pt] at (14,82.0)
    {\rotatebox{90}{migration transfer duration ($B/R$)}};

  % the floors each way is bound by, doubling as the gridlines
  \draw[grid] (32.0,70.5) -- (172.0,70.5);
  \node[ticklbl,anchor=north east] at (171.0,71.5) {$\tfrac{n-1}{2}B/R$};
  \draw[grid] (32.0,103.3) -- (172.0,103.3);
  \node[ticklbl,anchor=north east] at (171.0,104.3) {$B/R$};

  % axes
  \draw[axis] (32.0,136.0) -- (172.0,136.0);
  \draw[axis] (32.0,136.0) -- (32.0,28.0);
  \draw[axis] (32.0,136.0) -- (32.0,139.0);
  \node[ticklbl,anchor=north] at (32.0,140.0) {1};
  \draw[axis] (50.9,136.0) -- (50.9,139.0);
  \node[ticklbl,anchor=north] at (50.9,140.0) {2};
  \draw[axis] (88.6,136.0) -- (88.6,139.0);
  \node[ticklbl,anchor=north] at (88.6,140.0) {4};
  \draw[axis] (126.3,136.0) -- (126.3,139.0);
  \node[ticklbl,anchor=north] at (126.3,140.0) {6};
  \draw[axis] (164.0,136.0) -- (164.0,139.0);
  \node[ticklbl,anchor=north] at (164.0,140.0) {8};
  \draw[axis] (32.0,103.3) -- (29.0,103.3);
  \node[ticklbl,anchor=east] at (27.0,103.3) {1};
  \draw[axis] (32.0,70.5) -- (29.0,70.5);
  \node[ticklbl,anchor=east] at (27.0,70.5) {2};
  \draw[axis] (32.0,37.8) -- (29.0,37.8);
  \node[ticklbl,anchor=east] at (27.0,37.8) {3};
  \node[ticklbl,anchor=north] at (98.0,152.0) {chunks $k$};

  \node[ticklbl,anchor=north west] at (38,112) {$n = 5$};

  % (a) layered: k does not help, (1+f) B/R
  \draw[layered] (32.0,37.8) -- (164.0,37.8);
  \node[ticklbl,anchor=south east] at (164.0,36.3) {(a) layered};

  % (b) chain: (1 + f/k) B/R
  \draw[chain] (32.0,37.8) -- (50.9,70.5) -- (69.7,81.5) -- (88.6,86.9) -- (107.4,90.2) -- (126.3,92.4) -- (145.1,93.9) -- (164.0,95.1);
  \fill[ink] (32.0,37.8) circle (1.4pt);
  \fill[ink] (50.9,70.5) circle (1.4pt);
  \fill[ink] (69.7,81.5) circle (1.4pt);
  \fill[ink] (88.6,86.9) circle (1.4pt);
  \fill[ink] (107.4,90.2) circle (1.4pt);
  \fill[ink] (126.3,92.4) circle (1.4pt);
  \fill[ink] (145.1,93.9) circle (1.4pt);
  \fill[ink] (164.0,95.1) circle (1.4pt);
  \node[ticklbl,anchor=south] at (107.4,86.2) {(b) chain};
\end{tikzpicture}}%
}

\newcommand{\topolegend}{%
\scalebox{\topozoom}{%
\begin{tikzpicture}[x=1pt,y=-1pt,topostyle]
  \def\key##1{\tikz[baseline=-2.7pt]\draw[##1] (0pt,0pt) -- (18pt,0pt);\hspace{4pt}}
  \node[inner sep=0pt,anchor=base] at (0,0) {%
    \key{notify}Replicated log entry\hspace{18pt}%
    \key{bulk}Synchronous RDMA write\hspace{18pt}%
    \key{abulk}Asynchronous RDMA write};
\end{tikzpicture}}%
}

\begin{figure*}[t]
  \centering
  \begin{subfigure}[b]{0.32\textwidth}
    \centering\topopanel{2}
    \caption{migration layered on replication}
    \label{fig:topo-layered}
  \end{subfigure}\hfill
  \begin{subfigure}[b]{0.32\textwidth}
    \centering\topopanel{3}
    \caption{\OurImplementation{}, the integrated chain}
    \label{fig:topo-chain}
  \end{subfigure}\hfill
  \begin{subfigure}[b]{0.32\textwidth}
    \centering\topoplot
    \caption{what the split buys, by the model}
    \label{fig:topo-model}
  \end{subfigure}
  \par\vspace{4pt}\topolegend
  \caption{Two ways of landing a migrated group
  \textit{G}\textsubscript{\textit{i}} in a recipient shard group of $n = 5$ replicas.
  In both the donor leader A5 writes the blocks of
  \textit{G}\textsubscript{\textit{i}} straight into the recipient leader B5 with a
  one-sided RDMA write, and each leader drives its own \emph{replication path}, the
  leader-to-follower log that carries the control actions of the migration.
  What differs is how the data reaches the followers.
  In (a) the recipient leader sends every chunk itself, once per follower, so the cost of
  the transfer is paid again for each replica.
  In (b) each replica forwards to its successor, so the group becomes durable once a
  majority of $(n+1)/2 = 3$ replicas hold it, the shaded ones, and the remaining
  $(n-1)/2 = 2$ are served asynchronously by the rest of the chain.
  The time each one takes to make \textit{G}\textsubscript{\textit{i}} durable follows
  from the drawing, with $B$ the size of the group, $R$ the rate of one link, and $k$ the
  number of chunks the group is split into.
  In (a) the replication of the group follows its transfer rather than overlapping it, so
  splitting the group buys nothing and the majority costs $\frac{n+1}{2}B/R$ whatever
  $k$ is.
  In (b) every link carries $B$ once and the $(n+1)/2$ hops pipeline, so the majority
  costs $\bigl(1 + \frac{n-1}{2k}\bigr)B/R$ and falls towards $B/R$ as the group is split
  more finely.
  Panel (c) plots the two against $k$ at $n = 5$.
  They meet at $k = 1$, where the chain is as deep as the group is long and buys nothing,
  and the chain falls away as the group is split.
  The floor at $\frac{n-1}{2}B/R$ is the best a layered way could do even if it did
  overlap the two,
  because the recipient leader's own uplink still carries $B$ once per follower, and it
  remains above the $B/R$ the chain approaches.
  Panel (c) is the model rather than measured data.}
  \label{fig:topology}
\end{figure*}

\paragraph{Production deployments are replicated.}
% [Notes: every shard is a shard group of 2f+1 replicas with one leader, kept consistent by
%  a leader-based protocol (Raft), for durability and availability. Cite raft, paxos-simple,
%  slicer, shard-manager. State the consequence that frames the rest of the section: a
%  migration does not move data between two nodes but between two shard groups, so every
%  record that moves must also reach the followers of the recipient group before the
%  recipient can own it. Point the reader at Fig.~\ref{fig:topology}.]
\textcolor{red}{TBD}

\paragraph{A migration crosses two paths, not one.}
% [Notes: read off Fig.~\ref{fig:topology}. The replication path carries the control
%  actions of the migration and is a leader-to-follower fan-out with a majority
%  requirement. The chain-replication path carries the bulk data and is a pipeline whose
%  depth grows with the size of the group. The two have different shapes and different
%  costs, which is why a design that pushes the bulk data through the log, or that treats
%  the two as independent mechanisms, pays for it. Note what the figure shows about group
%  width: at five replicas the majority is reached after two chain hops rather than one.]
\textcolor{red}{TBD}

\paragraph{Layering migration on replication is expensive.}
% [Notes: this is the quantitative hook. Rocksteady~\cite{rocksteady} pulls records on
%  demand and defers re-replication until after the migration: on YCSB-B it falls from
%  452.0 to 124.8Kops/s, a loss of 72%, and spends 88s below its pre-migration level, with
%  the median read going from 22 to 45us and the 99.9th percentile from 64 to 255us. On
%  YCSB-A it falls from 232.0 to 144.0Kops/s, 38%, for 38s. Contrast with \OurImplementation{}
%  + FT, which gives up 0.9% over the seven seconds of the migration on YCSB-B and shows no
%  measurable loss on YCSB-A. Forward-reference Section~\ref{sec:eval-sota}.]
\textcolor{red}{TBD}

\paragraph{Integration makes replication cheap.}
% [Notes: the positive half of the argument. Because the replication of one group overlaps
%  the transfer of the next instead of following it, the four groups contribute 2.22s of
%  chain replication without extending the migration at all: the logged migration is
%  7.02--7.32s in every run, replicated or not. The price is paid in steady state instead,
%  6.4% of throughput on YCSB-B and 9.3% on YCSB-A, and in the update latency, which rises
%  from 5.49 to 9.72ms on YCSB-B, because an update is acknowledged only once a majority
%  holds it. Reads are unaffected, being served by the leader alone. Widening the group to
%  the five replicas of Fig.~\ref{fig:topology} adds one synchronous hop and costs 0.77s
%  over four groups, 2.22s against 2.99s. Forward-reference Section~\ref{sec:eval-replicas}.]
\textcolor{red}{TBD}

\paragraph{A migration in flight must survive a crash.}
% [Notes: the second gap, the fault-tolerance one. A migration spans two shard groups and
%  takes seconds, so it overlaps the failures the groups are already provisioned to
%  survive, yet prior work leaves the migration itself outside the guarantees of the
%  replication protocol: a crash of a donor or recipient leader abandons the transfer. State
%  what the section's argument demands of a design, which is what Section~\ref{sec:protocol}
%  then builds: control actions on the replicated log, bulk data off it, and recovery that
%  rolls a migration forward rather than restarting it. Measured cost: recovery in under
%  3.4s and completion less than 3s slower than a failure-free migration.]
\textcolor{red}{TBD}
